% Chapter 3, Topic _Linear Algebra_ Jim Hefferon
%  http://joshua.smcvt.edu/linearalgebra
%  2001-Jun-12
\topic{Geometri Pemetaan Linear} 
\index{geometri pemetaan linear|(}
%The geometry of linear maps %$\map{h}{\Re^n}{\Re^m}$ 
%is appealing both for its simplicity and for its usefulness.
Pasangan-pasangan gambar berikut membandingkan aksi geometris pemetaan nonlinear
\( f_1(x)=e^x \) dan \( f_2(x)=x^2 \)
\begin{center}
  \includegraphics{map/mp/ch3.46}
  \hspace*{4em}
  \includegraphics{map/mp/ch3.47}
\end{center}
dengan aksi pemetaan linear
\( h_1(x)=2x \) dan \( h_2(x)=-x \).
\begin{center}
  \includegraphics{map/mp/ch3.48}
  \hspace*{4em}
  \includegraphics{map/mp/ch3.49}
\end{center}
Keempat gambar tersebut menunjukkan domain $\Re$ di sebelah kiri yang dipetakan
ke kodomain $\Re$ di sebelah kanan.
Anak panah menelusuri ke mana setiap pemetaan mengirim
$x=0$, $x=1$, $x=2$, $x=-1$, dan $x=-2$.

Pemetaan nonlinear mendistorsi domain ketika mengubahnya menjadi ruang hasil.
Sebagai contoh,
jarak \( f_1(1) \) dari $f_1(2)$ lebih jauh daripada jaraknya dari $f_1(0)$
\Dash pemetaan ini meregangkan domain
secara tidak merata, sehingga interval domain di dekat $x=2$ diregangkan lebih
jauh daripada interval domain di dekat $x=0$.
Pemetaan linear lebih teratur: bagi setiap pemetaan, seluruh domain diregangkan
dengan faktor yang sama.
Pemetaan~$h_1$ di sebelah kiri membuat setiap interval dua kali lebih lebar,
sedangkan~$h_2$ di sebelah kanan mempertahankan panjang interval tetapi membalik
orientasinya, seperti interval naik dari $1$ ke $2$ yang diubah menjadi interval
turun dari $-1$ ke~$-2$.

Satu-satunya pemetaan linear dari $\Re$ ke $\Re$ adalah perkalian dengan skalar,
tetapi di dimensi yang lebih tinggi dapat terjadi lebih banyak hal.
Sebagai contoh, transformasi linear pada $\Re^2$ berikut memutar vektor
berlawanan arah jarum jam.\index{rotasi}\index{pemetaan linear!rotasi}
\begin{center}
  \includegraphics{map/mp/ch3.50}
\end{center}
Transformasi pada $\Re^3$ yang memproyeksikan vektor ke bidang $xz$ juga bukan
sekadar penskalaan ulang.
\begin{center}
 \includegraphics{map/mp/ch3.51}
\end{center}

Meskipun lebih beragam, pemetaan linear tetap berperilaku teratur bahkan di
dimensi yang lebih tinggi.
Tinjau pemetaan linear $\map{h}{\Re^n}{\Re^m}$ dan gunakan basis-basis standar
untuk merepresentasikannya dengan matriks $H$.
Ingat dari Teorema~V.\ref{th:CanonFormForMatEquiv} bahwa $H$ dapat difaktorkan
sebagai $H=PBQ$, dengan $P$ dan $Q$ nonsingular serta $B$ matriks identitas
parsial. Ingat pula bahwa matriks nonsingular dapat difaktorkan menjadi
matriks-matriks elementer\index{matriks!reduksi elementer}%
\index{matriks reduksi elementer}
$PBQ=T_nT_{n-1}\cdots T_sBT_{s-1}\cdots T_1$,
yaitu matriks yang diperoleh dari identitas $I$ setelah satu operasi baris
Gauss, sehingga setiap matriks $T$ memiliki salah satu dari tiga bentuk berikut:
\begin{equation*}
  I\grstep{k\rho_i}M_i(k) 
  \qquad 
  I\grstep{\rho_i\leftrightarrow\rho_j}P_{i,j}  
  \qquad
  I\grstep{k\rho_i+\rho_j}C_{i,j}(k) 
\end{equation*}
dengan $i\neq j$ dan $k\neq 0$.
Jadi, jika kita memahami efek geometris pemetaan linear yang dideskripsikan oleh
matriks identitas parsial dan efek pemetaan linear yang dideskripsikan oleh
matriks-matriks elementer, dalam arti tertentu kita akan memahami sepenuhnya
efek setiap pemetaan linear.
(Agar mudah digambar, gambar-gambar di bawah dibatasi pada transformasi $\Re^2$,
tetapi prinsipnya berlaku bagi pemetaan dari sebarang $\Re^n$ ke sebarang
$\Re^m$.)

Efek geometris transformasi linear yang direpresentasikan oleh matriks identitas
parsial adalah proyeksi.
\begin{equation*}
  \colvec{x \\ y  \\ z}
  \quad\xrightarrow{{\text{\tiny$\displaystyle\left(\begin{array}[r]{@{}c@{\hspace{0.8em}}c@{\hspace{0.8em}}c@{}}
                1  &0  &0 \\
                0  &1  &0 \\
                0  &0  &0
          \end{array}\right)$}}}
% \begin{mat}[r]
%       1  &0  &0   \\
%       0  &1  &0   \\
%       0  &0  &0   
%     \end{mat}} % _{\stdbasis_3,\stdbasis_3}}
  \qquad
  \colvec{x \\ y  \\ 0}
\end{equation*}

Efek geometris matriks $M_i(k)$ adalah meregangkan vektor dengan faktor $k$
sepanjang sumbu ke-$i$.
Pemetaan berikut meregangkan dengan faktor $3$ sepanjang sumbu $x$.
\begin{center}
  \includegraphics{map/mp/ch3.52}
\end{center}
Jika $0\leq k<1$ atau $k<0$, komponen ke-$i$ bergerak ke arah sebaliknya, dalam
gambar ini ke kiri.
\begin{center}
  \includegraphics{map/mp/ch3.53}
\end{center}
Kedua jenis peregangan ini disebut
\definend{dilatasi}.\index{dilatasi}\index{pemetaan linear!dilatasi}

Transformasi yang direpresentasikan oleh matriks $P_{i,j}$ menukar sumbu ke-$i$
dan ke-$j$. Transformasi ini merupakan \definend{pencerminan}%
\index{pencerminan}\index{pemetaan linear!pencerminan} terhadap hiperbidang
$x_i=x_j$.
\begin{center}
  \includegraphics{map/mp/ch3.54}
\end{center}
Permutasi yang melibatkan lebih dari dua sumbu dapat diuraikan menjadi kombinasi
pertukaran pasangan-pasangan sumbu; lihat
\nearbyexercise{exer:PermIsCompSwaps}.

Matriks-matriks yang tersisa memiliki bentuk $C_{i,j}(k)$.
Sebagai contoh, $C_{1,2}(2)$ menjalankan $2\rho_1+\rho_2$.
\begin{equation*}
  \colvec{x  \\  y}\quad
  \xrightarrow{\text{\tiny$\displaystyle\left(\begin{array}[r]{@{}c@{\hspace{0.8em}}c@{}}
                1  &0  \\
                2  &1  
          \end{array}\right)$}}  % _{\stdbasis_2,\stdbasis_2}}
   \quad\colvec{x \\ 2x+y}
\end{equation*}
Pada gambar di bawah, vektor $\vec{u}$ yang komponen pertamanya $1$ terkena efek
lebih kecil daripada vektor $\vec{v}$ yang komponen pertamanya $2$.
Vektor $\vec{u}$ dipetakan ke $h(\vec{u})$ yang hanya $2$ satuan lebih tinggi
daripada $\vec{u}$, sedangkan $h(\vec{v})$ berada $4$ satuan lebih tinggi
daripada $\vec{v}$.
\begin{center}
  \includegraphics{map/mp/ch3.55}
\end{center}
Setiap vektor yang komponen pertamanya $1$ akan terkena efek yang sama seperti
$\vec{u}$: vektor itu akan bergeser ke atas sejauh $2$.
Setiap vektor yang komponen pertamanya $2$ akan bergeser ke atas sejauh $4$,
seperti $\vec{v}$.
Artinya, transformasi yang direpresentasikan oleh $C_{i,j}(k)$ memengaruhi
vektor berdasarkan komponen ke-$i$-nya.

Cara lain untuk melihatnya adalah dengan meninjau aksi pemetaan ini pada persegi
satuan. Pada gambar berikut, vektor dengan komponen pertama $0$, seperti titik
asal, sama sekali tidak terdorong secara vertikal, sedangkan vektor dengan
komponen pertama positif bergeser ke atas.
Di sini semua vektor yang komponen pertamanya $1$, yakni seluruh sisi kanan
persegi, bergeser sejauh sama.
Secara umum, vektor-vektor pada garis vertikal yang sama bergeser sejauh sama,
yaitu dua kali komponen pertamanya.
Bangun yang dihasilkan memiliki alas dan tinggi yang sama dengan persegi
(sehingga luasnya juga sama), tetapi sudut-sudut siku-sikunya lenyap.
\begin{center}
  \includegraphics{map/mp/ch3.56}
\end{center}

Sebagai perbandingan, gambar berikut menunjukkan efek pemetaan yang
direpresentasikan oleh $C_{2,1}(2)$.
Di sini vektor dipengaruhi berdasarkan komponen keduanya:
$\binom{x}{y}$ bergeser secara horizontal sejauh dua kali $y$.
\begin{center}
  \includegraphics{map/mp/ch3.57}
\end{center}
Secara umum, bagi setiap $C_{i,j}(k)$, vektor-vektor dengan komponen ke-$i$ yang
sama bergeser sejauh sama.
Pemetaan semacam ini disebut
\definend{geseran}.\index{geseran}\index{pemetaan linear!geseran}

Dengan demikian, kita memahami efek geometris keempat jenis matriks pada ruas
kanan $H=T_nT_{n-1}\cdots T_sBT_{s-1}\cdots T_1$, sehingga dalam arti tertentu
kita memahami aksi sebarang matriks~$H$.
Jadi, bahkan di dimensi yang lebih tinggi, geometri pemetaan linear tetap
sederhana: geometri itu dibangun dengan merangkai sejumlah komponen yang
masing-masing beraksi secara sederhana.

Kita akan menerapkan pemahaman ini dalam dua cara.
Cara pertama adalah membuktikan suatu pernyataan umum tentang geometri pemetaan
linear. Ingat bahwa di bawah pemetaan linear, citra subruang merupakan subruang.
Karena itu, transformasi linear $h$ yang direpresentasikan oleh $H$ memetakan
garis yang melalui titik asal ke garis yang melalui titik asal.
(Dimensi ruang citra tidak dapat melebihi dimensi ruang domain, sehingga garis
tidak dapat dipetakan secara surjektif ke, misalnya, suatu bidang.)
Kita akan menunjukkan bahwa $h$ memetakan setiap garis\Dash bukan hanya garis
yang melalui titik asal\Dash ke suatu garis.
Pembuktiannya sederhana: proyeksi identitas parsial $B$ dan setiap $T_i$
elementer mengubah masukan berupa garis menjadi keluaran berupa garis;
verifikasi keempat kasus tersebut menjadi isi
\nearbyexercise{exer:ImageLinSurIsLinSur}.
Oleh karena itu, komposisinya juga mempertahankan garis.
% Thus, by understanding its components we can understand arbitrary square 
% matrices $H$, in the sense that we can prove things about them.

Cara kedua penerapan pemahaman geometris tentang pemetaan linear adalah untuk
memperjelas satu pokok dalam Kalkulus.
Gambar di bawah memperlihatkan aksi fungsi riil satu variabel
\( y(x)=x^2+x \).
Seperti fungsi-fungsi nonlinear yang digambarkan sebelumnya, efek geometris
pemetaan ini tidak teratur karena efeknya berbeda pada titik-titik domain yang
berbeda. Sebagai contoh, ketika masukan~$x$ bergerak dari $2$ ke $-2$, keluaran
terkait~$y(x)$ mula-mula menurun, berhenti sesaat, lalu meningkat.
\begin{center}
  \includegraphics{map/mp/ch3.58}
\end{center}
Namun, dalam Kalkulus kita tidak terlalu berfokus pada pemetaan secara
keseluruhan, melainkan pada efek lokalnya.
Di bawah ini kita mencermati perilaku pemetaan tersebut di dekat $x=1$.
Turunannya adalah $dy/dx=2x+1$, sehingga di dekat \( x=1 \) berlaku
\( \Delta y\approx 3\cdot\Delta x \).
%; in other words, \( (1.001^2+1.001)-(1^2+1)\approx 3\cdot (0.001) \).
Artinya, di suatu persekitaran $x=1$, pemetaan ini membuat domain membesar dengan
faktor $3$ \Dash secara lokal, pemetaan tersebut kira-kira merupakan dilatasi.
%The map $y$ is locally regular.
Gambar di bawah memperlihatkannya sebagai interval kecil
$(1-\Delta x\,..\,1+\Delta x)$ dalam domain yang dibawa ke interval
$(2-\Delta y\,..\,2+\Delta y)$ dalam kodomain yang tiga kali lebih lebar.
% $\Delta y \approx 3\cdot \Delta x$.
\begin{center}
  \includegraphics{map/mp/ch3.59}
\end{center}
% (When the above picture is drawn in the traditional cartesian way
% then the prior sentence about the rate of growth of $y(x)$ is usually
% stated:~the derivative $3$ is the slope of the 
% line tangent to the graph at the point $(1,2)$.)

Di dimensi yang lebih tinggi, gagasan intinya sama, tetapi lebih banyak hal dapat
terjadi. Bagi fungsi \( \map{y}{\Re^n}{\Re^m} \) dan titik
\( \vec{x}\in\Re^n \), turunan didefinisikan sebagai pemetaan linear
\( \map{h}{\Re^n}{\Re^m} \) yang paling baik menghampiri perubahan \( y \) di
dekat \( \vec{x} \).
Jadi, geometri yang dideskripsikan di atas berlaku langsung bagi turunan.
%Calculus considers the map 
%that locally approximates the change \( \Delta x\mapsto 3\cdot\Delta x \)
%(instead of the actual change map
%\( \Delta x\mapsto y(1+\Delta x)-y(1) \)) because the local map 
%is easier.
%It is easier in that, if the
%input change is doubled, or tripled, etc., then the
%resulting output change will double, or triple, etc.
%\begin{equation*}
%  3(r\,\Delta x)=r\,(3\Delta x)
%\end{equation*}
%(for $r\in\Re$), 
%and it is easier in that 
%adding changes in input adds the resulting output changes.
%\begin{equation*}
%  3(\Delta x_1+\Delta x_2)=3\Delta x_1+3\Delta x_2
%\end{equation*}
%In short, what's
%easy about \( \Delta x\mapsto 3\cdot\Delta x \) is that
%it is linear.

Sebagai penutup, kita perhatikan bagaimana sudut pandang ini memperjelas hasil
tentang turunan yang sering disalahpahami, yaitu Aturan Rantai.
Ingat bahwa, di bawah syarat-syarat yang sesuai pada kedua fungsi, turunan
komposisinya adalah sebagai berikut.
\begin{equation*}
  \frac{d\,(\composed{g}{f})}{dx}(x) = 
  \frac{dg}{dx}(f(x))\cdot\frac{df}{dx}(x)
\end{equation*} 
Sebagai contoh, turunan $\sin(x^2+3x)$ adalah
$\cos(x^2+3x)\cdot(2x+3)$.

Dari mana hasil ini berasal?
Tinjau $\map{f,g}{\Re}{\Re}$.
\begin{center}
  \includegraphics{map/mp/ch3.60}
\end{center}
Pemetaan pertama $f$ mendilatasi persekitaran $x$ dengan faktor
\begin{equation*}
  \frac{df}{dx}(x) 
\end{equation*}
kemudian pemetaan kedua $g$ mendilatasi persekitaran $f(x)$ dengan faktor
\begin{equation*}
  \frac{dg}{dx}(\,f(x)\,) 
\end{equation*}
dan ketika keduanya digabungkan, komposisi tersebut mendilatasi dengan hasil
kali kedua faktor itu.
Di dimensi yang lebih tinggi, pemetaan yang menyatakan perubahan fungsi di dekat
suatu titik merupakan pemetaan linear dan direpresentasikan oleh matriks.
Aturan Rantai mengalikan matriks-matriks tersebut.

% Thus, the geometry of linear maps $\map{h}{\Re^n}{\Re^m}$ 
% is appealing both for its simplicity and for its usefulness.

\begin{exercises}
  \item 
    Gunakan penguraian $H=PBQ$ untuk menemukan kombinasi dilatasi, pencerminan,
    geseran, dan proyeksi yang menghasilkan pemetaan
    $\map{h}{\Re^3}{\Re^3}$ yang direpresentasikan terhadap basis-basis standar
    oleh matriks berikut.
    \begin{equation*}
      H=\begin{mat}[r]
        1  &2  &1  \\
        3  &6  &0  \\
        1  &2  &2
      \end{mat}
    \end{equation*}
    \begin{answer}
      Reduksi Gauss berikut
      \begin{multline*}
        \grstep[-\rho_1+\rho_3]{-3\rho_1+\rho_2}
        \begin{mat}[r]
          1  &2  &1  \\
          0  &0  &-3 \\
          0  &0  &1
        \end{mat}
        \grstep{(1/3)\rho_2+\rho_3}
        \begin{mat}[r]
          1  &2  &1  \\
          0  &0  &-3 \\
          0  &0  &0
        \end{mat}                                     \\
        \grstep{(-1/3)\rho_2}
        \begin{mat}[r]
          1  &2  &1  \\
          0  &0  &1 \\
          0  &0  &0
        \end{mat}
        \grstep{-\rho_2+\rho_1}
        \begin{mat}[r]
          1  &2  &0  \\
          0  &0  &1 \\
          0  &0  &0
        \end{mat}
      \end{multline*}
      menghasilkan bentuk eselon tereduksi matriks tersebut.
      Selanjutnya, dua operasi kolom berupa menambahkan $-2$ kali kolom pertama
      ke kolom kedua, kemudian menukar kolom kedua dan ketiga, menghasilkan
      identitas parsial berikut.
      \begin{equation*} 
        B=\begin{mat}[r]
          1  &0  &0  \\
          0  &1  &0  \\ 
          0  &0  &0
        \end{mat}
      \end{equation*}
      Semua langkah itu dapat dinyatakan dalam bentuk matriks sebagai berikut:~jika
      \begin{equation*}
        R=
        \begin{mat}[r]
          1  &-1    &0  \\
          0  &1     &0  \\
          0  &0     &1         
        \end{mat}
        \begin{mat}[r]
          1  &0    &0  \\
          0  &-1/3 &0  \\
          0  &0    &1
        \end{mat}
        \begin{mat}[r]
          1  &0    &0  \\
          0  &1    &0  \\
          0  &1/3  &1         
        \end{mat}
        \begin{mat}[r]
          1  &0  &0  \\
          0  &1  &0  \\
         -1  &0  &1         
        \end{mat}
        \begin{mat}[r]
          1  &0  &0  \\
         -3  &1  &0  \\
          0  &0  &1         
        \end{mat}
      \end{equation*}
      dan
      \begin{equation*}
        C=
        \begin{mat}[r]
          1  &-2    &0  \\
          0  &1     &0  \\
          0  &0     &1         
        \end{mat}
        \begin{mat}[r]
          1  &0     &0  \\
          0  &0     &1  \\
          0  &1     &0         
        \end{mat}
      \end{equation*}
      maka $RHC=B$. Karena itu,
      $H=R^{-1}BC^{-1}$, yang merupakan faktorisasi berbentuk $PBQ$.
    \end{answer}
  \item 
    Kombinasi dilatasi, pencerminan, geseran, dan proyeksi apa yang menghasilkan
    rotasi sebesar $2\pi/3$ radian berlawanan arah jarum jam?
    \begin{answer}
      Pertama-tama, kita merepresentasikan pemetaan tersebut dengan matriks $H$,
      lalu menjalankan operasi baris dan, jika diperlukan, operasi kolom untuk
      mereduksinya menjadi matriks identitas parsial.
      Setelah itu, kita menerjemahkannya menjadi faktorisasi $H=PBQ$.
      Substitusi ke dalam matriks umum
          \begin{equation*}
            \rep{r_\theta}{\stdbasis_2,\stdbasis_2}
            =
            \begin{mat}
              \cos\theta  &-\sin\theta  \\
              \sin\theta  &\cos\theta
            \end{mat}
          \end{equation*}
          memberikan representasi berikut.
          \begin{equation*}
            \rep{r_{2\pi/3}}{\stdbasis_2,\stdbasis_2}
            =
            \begin{mat}[r]
              -1/2        &-\sqrt{3}/2  \\
              \sqrt{3}/2  &-1/2
            \end{mat}
          \end{equation*}
          Metode Gauss berjalan secara rutin.
          \begin{equation*}
            \grstep{\sqrt{3}\rho_1+\rho_2}
            \begin{mat}[r]
              -1/2        &-\sqrt{3}/2  \\
               0          &-2
            \end{mat}
            \grstep[(-1/2)\rho_2]{-2\rho_1}
            \begin{mat}[r]
               1          &\sqrt{3}    \\
               0          &1
            \end{mat}
            \grstep{-\sqrt{3}\rho_2+\rho_1}
            \begin{mat}[r]
               1          &0   \\
               0          &1
            \end{mat}
          \end{equation*}
          Hal itu dapat diterjemahkan menjadi persamaan matriks berikut.
          \begin{equation*}
            \begin{mat}[r]
              1  &-\sqrt{3}  \\
              0  &1
            \end{mat}
            \begin{mat}[r]
              -2  &0    \\
               0  &-1/2
            \end{mat}
            \begin{mat}[r]
               1         &0  \\
               \sqrt{3}  &1
            \end{mat}
            \begin{mat}[r]
              -1/2        &-\sqrt{3}/2  \\
              \sqrt{3}/2  &-1/2
            \end{mat}
            =I
          \end{equation*}
          Mengambil invers untuk menyelesaikan $H$ menghasilkan faktorisasi ini.
          \begin{equation*}
            \begin{mat}[r]
              -1/2        &-\sqrt{3}/2  \\
              \sqrt{3}/2  &-1/2
            \end{mat}
            =
            \begin{mat}[r]
                1         &0  \\
               -\sqrt{3}  &1
            \end{mat}
            \begin{mat}[r]
              -1/2  &0    \\
               0    &-2
            \end{mat}
            \begin{mat}[r]
              1  &\sqrt{3}  \\
              0  &1
            \end{mat}
            I
          \end{equation*}
    \end{answer}
  \item 
    Jika suatu pemetaan nonsingular, kita tidak memerlukan operasi kolom untuk
    mengubah representasinya menjadi matriks identitas, sehingga dalam $H=PBQ$
    matriks $Q$ merupakan identitas.
    Salah satu contoh pemetaan nonsingular adalah transformasi
    $\map{t_{-\pi/4}}{\Re^2}{\Re^2}$ yang memutar vektor sebesar $\pi/4$~radian
    searah jarum jam.
    \begin{exparts}
      \partsitem Temukan matriks $H$ yang merepresentasikan pemetaan ini terhadap
        basis-basis standar.
      \partsitem Gunakan Gauss--Jordan untuk mereduksi $H$ menjadi identitas
        tanpa operasi kolom.
      \partsitem
        Terjemahkan hasil tersebut menjadi persamaan matriks
        $T_jT_{j-1}\cdots T_1H=I$.
      \partsitem Selesaikan persamaan matriks itu untuk $H$.
      \partsitem Deskripsikan $H$ sebagai kombinasi dilatasi, pencerminan,
        geseran, dan proyeksi (identitas merupakan proyeksi trivial).
    \end{exparts}
    \begin{answer}
      \begin{exparts}
        \partsitem Ingat bahwa rotasi sebesar $\theta$~radian berlawanan arah
          jarum jam direpresentasikan terhadap basis standar sebagai berikut.
          \begin{equation*}
            \rep{t_{\theta}}{\stdbasis_2,\stdbasis_2}
            =\begin{mat}
              \cos\theta  &-\sin\theta  \\
              \sin\theta  &\cos\theta
             \end{mat}
          \end{equation*}
          Sudut searah jarum jam merupakan negatif dari sudut berlawanan arah
          jarum jam.
          \begin{equation*}
            \rep{t_{-\pi/4}}{\stdbasis_2,\stdbasis_2}
            =\begin{mat}
              \cos(-\pi/4)  &-\sin(-\pi/4)  \\
              \sin(-\pi/4)  &\cos(-\pi/4)
            \end{mat}
            =\begin{mat}[r]
              \sqrt{2}/2  &\sqrt{2}/2  \\
              -\sqrt{2}/2 &\sqrt{2}/2
            \end{mat}
          \end{equation*}
        \partsitem
          Reduksi Gauss--Jordan berikut
          \begin{equation*}
            \grstep{\rho_1+\rho_2}
            \begin{mat}[r]
              \sqrt{2}/2  &\sqrt{2}/2  \\
              0           &\sqrt{2}
            \end{mat}
            \grstep[(1/\sqrt{2})\rho_2]{(2/\sqrt{2})\rho_1}
            \begin{mat}[r]
              1  &1  \\
              0  &1
            \end{mat}
            \grstep{-\rho_2+\rho_1}
            \begin{mat}[r]
              1  &0  \\
              0  &1
            \end{mat}
          \end{equation*}
          menghasilkan matriks identitas.
          Jadi, kita tidak memerlukan operasi pertukaran kolom untuk berakhir
          pada identitas parsial.
        \partsitem Dalam perkalian matriks, reduksinya adalah
          \begin{equation*}
            \begin{mat}[r]
              1  &-1 \\
              0  &1
            \end{mat}
            \begin{mat}[r]
              2/\sqrt{2}  &0         \\
              0           &1/\sqrt{2}
            \end{mat}
            \begin{mat}[r]
              1  &0 \\
              1  &1
            \end{mat}
            H
            =I
          \end{equation*}
          (perhatikan bahwa komposisi operasi Gauss dibaca dari kanan ke kiri).
        \partsitem Mengambil invers
          \begin{equation*}
            H
            =
            \underbrace{
              \begin{mat}[r]
                1  &0 \\
                -1  &1
              \end{mat}
              \begin{mat}[r]
                \sqrt{2}/2  &0         \\
                0           &\sqrt{2}
              \end{mat}
              \begin{mat}[r]
                1  &1 \\
                0  &1
              \end{mat}
             }_P
            I
          \end{equation*}
          memberikan faktorisasi $H$ yang diinginkan.
          Identitas parsialnya adalah $I$.
          % , and $Q$ is trivial, that is, it is also an identity
          % matrix.
        \partsitem Dengan membaca komposisi dari kanan ke kiri (dan mengabaikan
          matriks identitas sebagai langkah trivial), kita memperoleh bahwa $H$
          berefek sama seperti pertama-tama menjalankan geseran berikut,
          \begin{center}
            \includegraphics{map/mp/ch3.95}
         \end{center}
         kemudian dilatasi yang mengalikan semua komponen pertama dengan
         $\sqrt{2}/2$ (ini merupakan penyusutan karena
         $\sqrt{2}/2\approx0.707$ kurang dari $1$) dan semua komponen kedua
         dengan $\sqrt{2}$, lalu geseran lain.
          \begin{center}
            \includegraphics{map/mp/ch3.96}
         \end{center}
         Sebagai contoh, kita mulai dengan vektor satuan yang sudutnya terhadap
         sumbu $x$ adalah $\pi/6$, lalu menerapkan komponen-komponen $H$ secara
         berurutan.
          \begin{center}
            \includegraphics{map/mp/ch3.97}
         \end{center}
         Kita dapat dengan mudah memverifikasi bahwa vektor yang dihasilkan
         memiliki panjang satu dan membentuk sudut $-\pi/12$ terhadap sumbu
         $x$, yang memang merupakan rotasi sebesar $\pi/4$ radian searah jarum
         jam karena $(\pi/6)-(\pi/4)=-\pi/12$.
      \end{exparts}
    \end{answer}
  \item \label{exer:RToRIsScalMult} 
    Tunjukkan bahwa setiap transformasi linear pada $\Re^1$ merupakan pemetaan
    $h_k$ yang mengalikan dengan skalar, $x\mapsto kx$.
    \begin{answer}
      Representasikan pemetaan tersebut terhadap basis-basis standar
      $\stdbasis_1,\stdbasis_1$.
      Hasilnya adalah matriks $\nbyn{1}$ yang satu-satunya entri merupakan
      skalar~$k$.
    \end{answer}
  \item \label{exer:ImageLinSurIsLinSur} 
    Tunjukkan bahwa pemetaan linear mempertahankan struktur linear suatu ruang.
    \begin{exparts}
      \partsitem Tunjukkan bahwa bagi setiap pemetaan linear dari $\Re^n$ ke
         $\Re^m$, citra setiap garis merupakan garis.
         Citranya dapat berupa garis degenerat, yaitu satu titik saja.
      \partsitem Tunjukkan bahwa citra setiap permukaan linear merupakan
         permukaan linear. Hasil ini memperumum fakta bahwa di bawah pemetaan
         linear, citra subruang merupakan subruang.
      \partsitem Pemetaan linear juga mempertahankan gagasan linear lainnya.
         Tunjukkan bahwa pemetaan linear mempertahankan sifat ``berada di
         antara'':~jika titik $B$ berada di antara $A$ dan $C$, maka citra $B$
         berada di antara citra $A$ dan citra $C$.
    \end{exparts}
    \begin{answer}
      \begin{exparts}
        \partsitem Garis merupakan himpunan bagian dari $\Re^n$ yang berbentuk
          $\set{\vec{v}=\vec{u}+t\cdot\vec{w}\suchthat t\in\Re}$.
          Citra titik pada garis tersebut adalah
          $h(\vec{v})=h(\vec{u}+t\cdot\vec{w})=h(\vec{u})+t\cdot h(\vec{w})$,
          dan himpunan vektor semacam itu, ketika $t$ menjelajahi bilangan riil,
          merupakan garis (meskipun degenerat jika $h(\vec{w})=\zero$).
        \partsitem Pernyataan ini merupakan perluasan langsung dari argumen
          sebelumnya.
        \partsitem Jika titik~$B$ berada di antara titik~$A$ dan~$C$, maka $B$
          terletak pada garis dari $A$ ke $C$.
          Artinya, terdapat $t\in (0\,..\,1)$ sedemikian sehingga
          $\vec{b}=\vec{a}+t\cdot (\vec{c}-\vec{a})$ (dengan $B$ sebagai titik
          ujung $\vec{b}$, dan seterusnya).
          Seperti dalam argumen butir pertama, linearitas menunjukkan bahwa
          $h(\vec{b})=h(\vec{a})+t\cdot h(\vec{c}-\vec{a})$.
      \end{exparts}
    \end{answer}
  \item 
    Gunakan gambar seperti yang muncul dalam pembahasan Aturan Rantai untuk
    menjawab pertanyaan berikut:~jika fungsi $\map{f}{\Re}{\Re}$ memiliki
    invers, apakah hubungan antara cara fungsi tersebut\Dash secara lokal dan
    hampiran\Dash mendilatasi ruang dan cara inversnya mendilatasi ruang?
    \begin{answer}
      Kedua faktor dilatasi itu saling invers.
      Sebagai contoh, bagi $x\in\Re$ yang tetap, jika $f^\prime (x)=k$ (dengan
      $k\neq 0$), maka $(f^{-1})^\prime (f(x))=1/k$.
      \begin{center}
        \includegraphics{map/mp/ch3.98}
     \end{center}
    \end{answer}
  \item \label{exer:PermIsCompSwaps}
    Tunjukkan bahwa bagi sebarang permutasi, yakni sebarang pengurutan ulang $p$
    atas bilangan $1$, \ldots, $n$, pemetaan
    \begin{equation*}
      \colvec{x_1 \\ x_2 \\ \vdots\\ x_n}
      \mapsto
      \colvec{x_{p(1)} \\ x_{p(2)} \\ \vdots \\ x_{p(n)}}
    \end{equation*}
    dapat dijalankan sebagai komposisi pemetaan-pemetaan yang masing-masing hanya
    menukar satu pasangan koordinat.
    \textit{Petunjuk:} Anda dapat menggunakan induksi pada $n$.
    (\textit{Catatan:}~dalam bab keempat kita akan menunjukkan hasil ini dan juga
    menunjukkan bahwa paritas banyak pertukaran yang digunakan ditentukan oleh
    $p$. Artinya, meskipun suatu permutasi dapat dinyatakan dengan dua cara yang
    menggunakan banyak pertukaran berbeda, kedua cara itu sama-sama menggunakan
    banyak pertukaran genap atau sama-sama menggunakan banyak pertukaran ganjil.)
    \begin{answer}
      Kita dapat menunjukkannya dengan induksi pada banyak komponen vektor.
      Dalam kasus dasar $n=1$, satu-satunya permutasi adalah permutasi trivial,
      dan pemetaan
      \begin{equation*}
        \colvec{x_1}
        \mapsto
        \colvec{x_1}
      \end{equation*}
      dapat dinyatakan sebagai komposisi pertukaran\Dash yaitu nol pertukaran.
      Untuk langkah induksi, kita mengasumsikan bahwa pemetaan yang diinduksi oleh
      sebarang permutasi atas kurang dari $n$ bilangan dapat dinyatakan hanya
      dengan pertukaran, lalu meninjau pemetaan yang diinduksi oleh permutasi $p$
      atas $n$ bilangan.
      \begin{equation*}
        \colvec{x_1 \\ x_2 \\ \vdots \\ x_n}
        \mapsto
        \colvec{x_{p(1)} \\ x_{p(2)} \\ \vdots \\ x_{p(n)}}
      \end{equation*}
      Tinjau bilangan~$i$ sedemikian sehingga $p(i)=n$.
      Pemetaan
      \begin{equation*}
        \colvec{x_1      \\ x_2      \\ \vdots \\ x_i      \\ \vdots \\ x_n}
        \mapsunder{\hat{p}}
        \colvec{x_{p(1)} \\ x_{p(2)} \\ \vdots \\ x_{p(n)} \\ \vdots  \\ x_{n}}
      \end{equation*}
      jika diikuti dengan pertukaran komponen ke-$i$ dan ke-$n$, akan menghasilkan
      pemetaan~$p$.
      Hipotesis induksi menyatakan bahwa $\hat{p}$ dapat diwujudkan sebagai
      komposisi pertukaran.
    \end{answer}
\end{exercises}
\index{geometri pemetaan linear|)}
\endinput
